\documentclass[10pt,a4paper]{amsart}
\usepackage[margin=19mm]{geometry}
\usepackage{amsmath,amssymb,mathtools}
\usepackage[scr=rsfs,cal=euler]{mathalfa}
\usepackage{xfrac}
\usepackage[unicode,hidelinks,bookmarks=false]{hyperref}
\newcommand{\ie}{i.e.\ }
\newcommand{\cf}{cf.\ }
\newcommand{\resp}{respectively\ }
\newcommand{\Subsection}{\subsection}
\providecommand{\nobreakdash}{\nobreak\hbox{-}}
\setlength{\emergencystretch}{2em}
\multlinegap0pt
\usepackage{amssymb}
\usepackage{enumerate}
\usepackage{xcolor}
\usepackage{mathtools}
\usepackage{tikz}
\usepackage{dsfont}
\usepackage{bm}
\newtheorem{lemma}{Lemma}
\title{Drinfeld Isomorphism for Novel Quantum Affine Algebra of Type \(A_{1}^{(1)}\)}
\author{Rushu Zhuang, Ge Feng, Naihong Hu$^*$}
\date{Source: arXiv:2601.20562v1 (CC BY 4.0)}
\begin{document}
\maketitle

\noindent\emph{Source and licence.} Drinfeld Isomorphism for Novel Quantum Affine Algebra of Type \(A_{1}^{(1)}\). Version arXiv:2601.20562v1 (CC BY 4.0), distributed by arXiv under the Creative Commons Attribution 4.0 International licence, http://creativecommons.org/licenses/by/4.0/, as recorded in the arXiv licence metadata for this exact version and shown on the abstract page https://arxiv.org/abs/2601.20562v1. Full licence notice: \texttt{licenses/arxiv-2601.20562-CC-BY-4.0.txt}.\par\smallskip

\noindent\textit{Editorial scope.} Counted unit: the article's lemma lemma (source0.tex lines 435--440). The article's complete original proof of it is delivered with it (source0.tex lines 441--523, one \texttt{proof} environment of 83 lines). No mathematics is written, repaired or re-derived: the statement and the proof are the article's own lines, in the article's own notation, and no retained line is substituted or re-flowed.  No further source-local context is needed: every cross-reference of every retained line resolves inside the two delivered spans.  Nothing was omitted from any retained span, so the package carries no omission record.  No retained line contains a citation, so the wrapper carries no bibliography and no bibliography entry.  No retained line contains a figure environment, an image-inclusion command or drawing code, so no image is delivered and the package declares no asset dependency.  Editorial formatting only, in addition: an AMS \texttt{amsart} preamble that restates the article's own declarations and macros used by the retained lines, namely the environments \texttt{lemma} on separate counters, together with the standard packages those lines need.  The retained environments print in this document as Lemma 1 in order of appearance, whereas the article prints the same items with the numbering of its own full document.  The complete native source of the article as distributed in the arXiv e-print for this exact version is bundled unchanged as \texttt{sources/193529/source0.tex}.
\begin{lemma}
For \(n,m\in\mathbb{Z}_{\ge 0}\), \(n\ge m\), the commutation relations between \(\mathcal{E}_{n\delta+\alpha_{1}}\) and \(\mathcal{E}_{m\delta+\alpha_{1}}\) are 
\begin{align*}     
\mathcal{E}_{n\delta+\alpha_{1}}&\mathcal{E}_{m\delta+\alpha_{1}}+(-1)^{n+m}q^{-2}\mathcal{E}_{m\delta+\alpha_{1}}\mathcal{E}_{n\delta+\alpha_{1}}\\ &=q^{-2}\mathcal{E}_{(n-1)\delta+\alpha_{1}}\mathcal{E}_{(m+1)\delta+\alpha_{1}}+(-1)^{n+m}\mathcal{E}_{(m+1)\delta+\alpha_{1}}\mathcal{E}_{(n-1)\delta+\alpha_{1}}.  
\end{align*}  
\end{lemma} 

\begin{proof}
We prove the lemma for indices \( n \) and \( m \) by double induction. First, we verify that the base case \( P(n,0) \) holds. Second, we assume that the lemma \( P(n,i) \) is true for all positive integers \( i \)  satisfying \( i<m \), then proceed to prove that \( P(n,m) \) holds.

(I) When \(m = 0\), let us consider the relations 
\begin{align}
\mathcal{E}_{n\delta+\alpha_1} e_1 + (-1)^n q^{-2} e_1 \mathcal{E}_{n\delta+\alpha_1} = q^{-2} \mathcal{E}_{(n-1)\delta+\alpha_1} \mathcal{E}_{\delta+\alpha_1} + (-1)^n \mathcal{E}_{\delta+\alpha_1} \mathcal{E}_{(n-1)\delta+\alpha_1}. 
\end{align}

We first examine the case of \(n=1\), then extend the result to general \(n \geq 1\).
 
Substituting \(n = 1\) and \(m = 0\) into (3.1), we need to prove 
\(
\mathcal{E}_{\delta+\alpha_1} e_1 =q^{-2} e_1 \mathcal{E}_{\delta+\alpha_1}.
\) It is now straightforward to verify that 
\[
\mathcal{E}_{\delta+\alpha_1} e_1 - q^{-2} e_1 \mathcal{E}_{\delta+\alpha_1}=\mathcal{T}_1\Phi\left(-e_1k^{-1}_2f_2 + q^{-2}k^{-1}_2f_2e_1\right).
\]    
We use \((A2)\) and \((A3)\) to get that
\(
\mathcal{E}_{\delta+\alpha_1} e_1 =q^{-2} e_1 \mathcal{E}_{\delta+\alpha_1}\).
 
Assume relation (3.1) holds for \(n = k\,(k>1)\). To confirm it holds for \(n = k+1\), we need to use the following identities
\begin{align}
[2]_q\mathcal{E}_{(t+1)\delta+\alpha_{1}} = \mathcal{E}_{\delta}\mathcal{E}_{t\delta+\alpha_{1}} + \mathcal{E}_{t\delta+\alpha_{1}}\mathcal{E}_{\delta},\quad \forall t\ge 0,\qquad (\text{by Theorem 2.10})
\end{align}
and use induction hypothesis to get
\begin{align}
\mathcal{E}_{s\delta+\alpha_1} e_1 + (-1)^s q^{-2} e_1 \mathcal{E}_{s\delta+\alpha_1} = q^{-2} \mathcal{E}_{(s-1)\delta+\alpha_1} \mathcal{E}_{\delta+\alpha_1} + (-1)^s \mathcal{E}_{\delta+\alpha_1} \mathcal{E}_{(s-1)\delta+\alpha_1},
\end{align}
where \(1\le s\le k\).

Applying (3.2) and (3.3), it is straightforward to verify that
\begin{align*}
[2]_q\mathcal{E}_{(k+1)\delta+\alpha_{1}}e_{1} &=\Bigl\{\mathcal{E}_{\delta}\mathcal{E}_{k\delta+\alpha_{1}} + \mathcal{E}_{k\delta+\alpha_{1}}\mathcal{E}_{\delta}\Bigr\}e_{1}\\
&=\mathcal{E}_{\delta}\Bigl\{(-1)^{k-1} q^{-2}e_1\mathcal{E}_{k\delta+\alpha_{1}}+q^{-2}\mathcal{E}_{(k-1)\delta+\alpha_{1}}\mathcal{E}_{\delta+\alpha_{1}}+(-1)^k\mathcal{E}_{\delta+\alpha_{1}}\mathcal{E}_{(k-1)\delta+\alpha_{1}}\Bigr\}\\
&\quad+\mathcal{E}_{k\delta+\alpha_{1}}\Bigl\{[2]_q\mathcal{E}_{\delta+\alpha_{1}} -e_1\mathcal{E}_{\delta}\Bigr\}\\
&=(-1)^{k-1} q^{-2}\Bigl\{[2]_q\mathcal{E}_{\delta+\alpha_{1}} -e_1\mathcal{E}_{\delta}\Bigr\}\mathcal{E}_{k\delta+\alpha_{1}}+q^{-2}\Bigl\{ [2]_q\mathcal{E}_{k\delta+\alpha_{1}}-\mathcal{E}_{(k-1)\delta+\alpha_{1}}\mathcal{E}_{\delta}\Bigr\}\mathcal{E}_{\delta+\alpha_{1}}
\\
&\quad+(-1)^k\Bigl\{[2]_q\mathcal{E}_{2\delta+\alpha_{1}}-\mathcal{E}_{\delta+\alpha_{1}}\mathcal{E}_{\delta}\Bigr\}\mathcal{E}_{(k-1)\delta+\alpha_{1}}+[2]_q\mathcal{E}_{k\delta+\alpha_{1}} \mathcal{E}_{\delta+\alpha_{1}}
\\
&\quad-\Bigl\{(-1)^{k-1} q^{-2}e_1\mathcal{E}_{k\delta+\alpha_{1}}+q^{-2}\mathcal{E}_{(k-1)\delta+\alpha_{1}}\mathcal{E}_{\delta+\alpha_{1}}+(-1)^k\mathcal{E}_{\delta+\alpha_{1}}\mathcal{E}_{(k-1)\delta+\alpha_{1}}\Bigr\}\mathcal{E}_{\delta}
\end{align*}
\begin{align*}
&=[2]_q\Bigl\{(-1)^{k}q^{-2}e_{1}\mathcal{E}_{(k+1)\delta+\alpha_{1}}+q^{-2}\mathcal{E}_{k\delta+\alpha_{1}}\mathcal{E}_{\delta+\alpha_{1}}+(-1)^{k+1}\mathcal{E}_{\delta+\alpha_{1}}\mathcal{E}_{k\delta+\alpha_{1}}\\
&\quad+\Bigl(\mathcal{E}_{k\delta+\alpha_{1}}\mathcal{E}_{\delta+\alpha_1}+(-1)^{k+1}q^{-2}\mathcal{E}_{\delta+\alpha_1}\mathcal{E}_{k\delta+\alpha_{1}}-q^{-2}\mathcal{E}_{(k-1)\delta+\alpha_{1}}\mathcal{E}_{2\delta+\alpha_{1}}\\
&\quad+(-1)^{k}\mathcal{E}_{2\delta+\alpha_{1}}\mathcal{E}_{(k-1)\delta+\alpha_{1}}\Bigr)\Bigr\}\\
&=[2]_q\Bigl\{(-1)^{k}q^{-2}e_{1}\mathcal{E}_{(k+1)\delta+\alpha_{1}}+q^{-2}\mathcal{E}_{k\delta+\alpha_{1}}\mathcal{E}_{\delta+\alpha_{1}}+(-1)^{k+1}\mathcal{E}_{\delta+\alpha_{1}}\mathcal{E}_{k\delta+\alpha_{1}}\\
&\quad+(\mathcal{T}_1\Phi)\Bigl(\mathcal{E}_{(k-1)\delta+\alpha_{1}}e_1+(-1)^{k+1}q^{-2}e_1\mathcal{E}_{(k-1)\delta+\alpha_{1}}-q^{-2}\mathcal{E}_{(k-2)\delta+\alpha_{1}}\mathcal{E}_{\delta+\alpha_{1}}\\ 
&\quad+(-1)^{k}\mathcal{E}_{\delta+\alpha_{1}}\mathcal{E}_{(k-2)\delta+\alpha_{1}}\Bigr)\Bigr\}.
\end{align*}
We ultimately obtain  
\begin{align*}
\mathcal{E}_{(k+1)\delta+\alpha_{1}}e_{1}+(-1)^{k+1}q^{-2}e_{1}\mathcal{E}_{(k+1)\delta+\alpha_{1}}
=q^{-2}\mathcal{E}_{k\delta+\alpha_{1}}\mathcal{E}_{\delta+\alpha_{1}}+(-1)^{k+1}\mathcal{E}_{\delta+\alpha_{1}}\mathcal{E}_{k\delta+\alpha_{1}}.
\end{align*}

Thus, relation (3.1) holds for all \(n > 0\).

(II) When \(m=t\), \(n>t>0\), we assume the lemma holds for \(m=t\):
\begin{align}
\mathcal{E}_{n\delta+\alpha_1} \mathcal{E}_{t\delta+\alpha_1} + (-1)^{n+t}& q^{-2} \mathcal{E}_{t\delta+\alpha_1} \mathcal{E}_{n\delta+\alpha_1} \nonumber\\
&= q^{-2} \mathcal{E}_{(n-1)\delta+\alpha_1} \mathcal{E}_{(t+1)\delta+\alpha_1} + (-1)^{n+t} \mathcal{E}_{(t+1)\delta+\alpha_1} \mathcal{E}_{(n-1)\delta+\alpha_1}.
\end{align}
We claim that the relation holds for \(m=t+1\).

To prove the claim, applying the automorphism \(\mathcal{T}_1\Phi\) to relation (3.4), we obtain
\begin{align}
\mathcal{E}_{(n+1)\delta+\alpha_1} \mathcal{E}_{(t+1)\delta+\alpha_1} + &(-1)^{n+t}q^{-2} \mathcal{E}_{(t+1)\delta+\alpha_1} \mathcal{E}_{(n+1)\delta+\alpha_1} \nonumber\\
&\quad = q^{-2} \mathcal{E}_{n\delta+\alpha_1}\mathcal{E}_{(t+2)\delta+\alpha_1} + (-1)^{n+t} \mathcal{E}_{(t+2)\delta+\alpha_1} \mathcal{E}_{n\delta+\alpha_1}.
\end{align}

We set \(n' = n+1\) and \(t' = t+1\). Since \(n > t\) implies \(n' > t'\), relation (3.5) simplifies to
\[
\begin{aligned}
\mathcal{E}_{n'\delta+\alpha_1} \mathcal{E}_{t'\delta+\alpha_1} + &(-1)^{n'+t'-1} q^{-2} \mathcal{E}_{t'\delta+\alpha_1} \mathcal{E}_{n'\delta+\alpha_1} \\
&\quad = q^{-2} \mathcal{E}_{(n'-1)\delta+\alpha_1} \mathcal{E}_{(t'+1)\delta+\alpha_1} + (-1)^{n'+t'-1} \mathcal{E}_{(t'+1)\delta+\alpha_1} \mathcal{E}_{(n'-1)\delta+\alpha_1},
\end{aligned}
\]
which matches the relation in the lemma for \(m = t' = t+1\), confirming the claim.

(III)
By the base case \(m = 0\), and inductive step, up to \(m = t+1\), the lemma holds.
\end{proof}

\end{document}
